\section{Introduction}

Autoregressive decoding in large language models (LLMs) is inherently
sequential: each generated token requires a full forward pass of the target
model. This constraint directly limits single-stream throughput in interactive
inference settings. Speculative decoding mitigates this bottleneck by adding a
smaller draft model that proposes a length-$k$ token block, followed by one
target-side verification step; modified rejection sampling preserves the target
distribution~\cite{Leviathan2023, Chen2023specsampling}.

Although speculative decoding is now well studied, deployment guidance remains
uncertain for consumer-GPU settings. Reported gains in the literature are often
obtained on data-center hardware and are sensitive to model pairing, proposal
length, and sampling regime. For practitioners operating on a single RTX-class
GPU, the key questions are practical: which draft size is worthwhile, which
$k$ is optimal, and whether deterministic versus stochastic decoding changes
the speed--quality frontier.

This study addresses three practical gaps for deployment-facing evidence:
(i) broader hardware/model validation, (ii) clearer seed-level variance and
statistical testing details, and (iii) explicit positioning of Adaptive
Cascaded Speculative Decoding (ACSD) as a robustness policy rather than a
speedup-maximizing algorithm.

This direction matters for privacy-sensitive and latency-critical local use,
including local workstations and edge-side decision-support workflows. We
therefore position this paper as a hardware-aware characterization study,
rather than a new decoding algorithm.

Within this scope, ACSD is not introduced as a mechanism to exceed the peak
throughput of the best static speculative setting. Instead, we position ACSD as
a hardware-aware tail-latency guardrail for dynamically difficult inputs,
trading a small amount of mean throughput for more controlled runtime behavior
under acceptance drops.

\subsection{Research questions}
\label{sec:rqs}

We address four research questions:
\begin{enumerate}\itemsep0pt
  \item[\textbf{RQ1.}] How much end-to-end speedup does vanilla
        speculative decoding deliver against autoregressive decoding on
      consumer RTX hardware with Qwen2.5 and Qwen3 family pairings, and how
      close is this to
        a practical hardware-constrained upper bound?
  \item[\textbf{RQ2.}] How do draft model size and draft length jointly
        determine token acceptance and net latency, and what do these trends
        imply about memory-traffic and kernel-launch overhead on consumer GPUs?
  \item[\textbf{RQ3.}] How sensitive is the speed--quality trade-off to
        the sampling regime (deterministic vs.\ stochastic) and to task
        entropy?
  \item[\textbf{RQ4.}] Do the ranking trends (optimal $k$, regime preference)
        transfer across consumer-GPU profiles (desktop RTX4090 vs.
        RTX5090-laptop), and where do absolute gains fail to transfer?
\end{enumerate}

\subsection{Contributions}
\label{sec:contributions}

This paper makes the following contributions:
\begin{enumerate}\itemsep0pt
  \item A reproducible consumer-GPU protocol with frozen sample manifests,
        matched prompts, and two decoding regimes across GSM8K, MMLU subset,
        and CNN/DailyMail.
  \item A two-family fixed-policy study combining Qwen2.5 and Qwen3
        configurations, with $k\in\{4,8,16\}$ and
        deterministic/stochastic regimes, reported with latency, TTFT, TPOT,
        acceptance, and quality metrics.
  \item Cross-hardware validation across desktop RTX4090 and RTX5090-laptop
        configurations, showing trend portability (best at low $k$,
        deterministic preference) but non-portable absolute gains.
  \\item Expanded statistical reporting: seed-level
        mean$\pm$std summaries, paired tests with Holm-Bonferroni correction,
        and explicit effect reporting.
  \\item An ACSD interpretation and ablation that prioritizes runtime
        stability and tail behavior over peak throughput claims.
\end{enumerate}

The rest of the paper is organized as follows. Section~\ref{sec:related}
reviews related work. Section~\ref{sec:protocol} defines the experimental
protocol; Section~\ref{sec:metrics} defines metrics. Section~\ref{sec:impl}
describes implementation and reproducibility controls. Section~\ref{sec:results}
reports the main empirical findings. Section~\ref{sec:acsd} outlines
an adaptive speculative-strategy framework and its observed runtime behavior.
Section~\ref{sec:discussion}
discusses implications, limitations, and conclusions.
